\section{The connection part of the shift}
\label{sec:bfield}

A constant shift of the interior Lagrangian has two summands (\cref{thm:EAdichotomy}): an $\eta$-antisymmetric one, in $\so{3,3}$, and an
$\eta$-self-adjoint one. \Cref{sec:neutral,sec:closure,sec:restclosure} were built on the second, the binding term $\kappa G$, and met the first only as transport along a path (\cref{prop:tear}) and in \cref{thm:EAdichotomy}. This section adds the
first, a constant $b\in\so{3,3}$, and asks what it does to the binding. Only $b$ is a connection, and where it is gauged it is the value of a
Yang--Mills potential along the worldline, the ``$b$-field'' in the notation of Yang and Mills \cite{YangMills54}. In brief, the answers are as follows. For
every $b$ the field equation factorises covariantly and the binding term transforms as a tensor, so no gauge transformation creates it
(\cref{bf:covariant}). An admissible gauge potential is compact. Its part commuting with $G$ splits the binding frequency into Zeeman lines that a
moving frame removes (\cref{bf:larmor,bf:lG}); a purely anticommuting part tears the binding (\cref{bf:pG}); a mixed one keeps every orbit bounded
exactly when the two triples rotate at equal rates about axes that are not opposite, outside one exceptional ratio (\cref{bf:screen}). On its own
a connection does not bind, and nothing in the interior dynamics ties the mass $m$ to the rate $\kappa$ (\cref{bf:alone}). Wherever the motion is
bounded the zero-point energy cancels frequency by frequency (\cref{bf:krein}). At three constituents the part of the Pati--Salam algebra
compatible with the binding splits the frequency by $\su2$ weights (\cref{bf:so4spec}). None of this supplies a scale. That the scale of a quantum
Yang--Mills theory is the binding length is stated as a conjecture (\cref{bf:transmute}).

Units are $\hbar=c=1$ unless a symbol is shown. The parameter $\tau$ has the dimension of a time, $\kappa=c/\ell$ of \eqref{eq:kappa} and every
eigenvalue of $b$ are rates, and $\hbar\kappa$ is a quantum of interior energy.

\paragraph{Notation.} $V=\R^{3,3}$, $F^\dagger:=\eta^{-1}F^{\!\top}\eta$ is the $\eta$-adjoint, $G$ is the binding root \eqref{eq:Gpair} and $\cK$
the grading root of \cref{prop:twoK}; on vectors $\cK=\tfrac12\ad\bigl(\sum_ie_ie_{i+3}\bigr)$ and $G=-\cK\eta$ (\cref{lem:Gclifford}). For indices $a<b$ (here $b$ is an index, not the connection) the
element $R_{ab}\in\so{3,3}$ is defined by $R_{ab}e_b=\eta_{bb}e_a$, $R_{ab}e_a=-\eta_{aa}e_b$, and zero on the other basis vectors; it is
$\tfrac12\ad(e_ae_b)$ on vectors. The generators $L_i,K_i$ of $\lG$ (\cref{thm:emfield}) act on $z=\bB+i\bE$ as $z\mapsto\mathbf e_i\times z$ and
$z\mapsto i\,\mathbf e_i\times z$, so $K_i=GL_i$, the $L_i$ are rotations and the $K_i$ boosts, and
\begin{equation}
\label{bf:eq:w}
b=\textstyle\sum_i\bigl(w^{\rm r}_iL_i+w^{\rm b}_iK_i\bigr)\in\lG\quad\text{acts as}\quad z\mapsto w\times z,\qquad w=w^{\rm r}+i\,w^{\rm b}\in\C^3 .
\end{equation}
In particular $R_{12}-R_{45}=-L_3$. We write $\sigma_j$ for the boost exchanging $e_j$ and $e_{j+3}$, $[\mathbf v]_\times$ for the matrix of
$\mathbf x\mapsto\mathbf v\times\mathbf x$ on $\R^3$, $\mathcal E_\mu(T)=\ker(T-\mu)$ for the eigenspaces of an operator on $V\otimes\C$, and
$h_\eta(y,y')=\bar y^{\!\top}\eta\,y'$ for the Hermitian extension of $\eta$. The interior Lagrangian with a constant shift is
\begin{equation}
\label{bf:eq:L}
\mathcal L=\frac m2\,\eta(u,u),\qquad u=\dot X+SX,\qquad S=b+\kappa G,\qquad b\in\so{3,3},
\end{equation}
with $p=m\,\eta u$ and $\pi=\eta^{-1}p=mu$. For $b=0$ it is the split-form Lagrangian \eqref{eq:splitL}.

\subsection{Two summands, and which one is a connection}
\label{sec:bf-split}

By \cref{thm:EAdichotomy} at $n=1$, $\End(V)=\so{3,3}\oplus\{F:F=F^\dagger\}$, of dimensions $15+21$, and by \cref{prop:twoK} the binding root
lies in the second summand and the grading root in the first. The $\eta$-antisymmetric part of $S$ is therefore $b$, and its self-adjoint part
is $\kappa G$. The decomposition is linear algebra, valid for any nondegenerate form; the next theorem gives it a dynamical meaning.

\begin{theorem}[Covariant factorisation]
\label{bf:covariant}
Let $b\in\so{3,3}$ be constant and $\nabla:=\frac{d}{d\tau}+b$.
\begin{enumerate}[label=(\roman*),leftmargin=2.2em]
\item The Euler--Lagrange equations of \eqref{bf:eq:L} are $(\nabla-\kappa G)\pi=0$ with $\pi=m(\nabla+\kappa G)X$, that is
\begin{equation}
\label{bf:eq:covEL}
(\nabla-\kappa G)(\nabla+\kappa G)X=0\qquad\Longleftrightarrow\qquad \nabla^2X+\kappa^2X+\kappa\,(\nabla_{\!\tau}G)\,X=0,\qquad \nabla_{\!\tau}G:=[b,G].
\end{equation}
\item Let $X=g(\tau)Y$ with $g(\tau)\in O(3,3)$. Then $\bigl(\tfrac d{d\tau}+S\bigr)(gY)=g\bigl(\tfrac d{d\tau}+S^g\bigr)Y$ with
$S^g=g^{-1}Sg+g^{-1}\dot g$. The $\so{3,3}$ part $b^g=g^{-1}bg+g^{-1}\dot g$ of $S^g$ transforms as a connection, its self-adjoint part
$\kappa\,g^{-1}Gg$ as a tensor, and $\tfrac d{d\tau}(g^{-1}Gg)+[b^g,g^{-1}Gg]=g^{-1}(\nabla_{\!\tau}G)g$.
\end{enumerate}
\end{theorem}

\begin{proof}
(i) $\partial\mathcal L/\partial\dot X=p$ and $\partial\mathcal L/\partial X=S^{\!\top}p$, so $\dot p=S^{\!\top}p$, that is
$\dot\pi=S^\dagger\pi=(-b+\kappa G)\pi$, which is $(\nabla-\kappa G)\pi=0$. Moreover $\pi=m(\nabla+\kappa G)X$, $\nabla(GX)-G\nabla X=[b,G]X$ and
$G^2=-\id$. (ii) Differentiating $g^{\!\top}\eta g=\eta$ gives $(g^{-1}\dot g)^\dagger=-g^{-1}\dot g$, and $F\mapsto g^{-1}Fg$ commutes with the
$\eta$-adjoint, so it preserves both summands; the last identity is a direct differentiation.
\end{proof}

Part~(ii) is inherited: it is the transformation law of a connection and of an endomorphism field in any gauge theory. Part~(i) is fixed. It is
the internal-connection, split-form version of the curved second-order equation at the end of \cref{sec:shifted}, which, as there, adds to the case $b=0$ only the term
$\kappa(\nabla_{\!\tau}G)X$. That the binding part transforms as a tensor restates $G=G^\dagger$: no gauge transformation creates, removes or
rescales it. The connection is the transport part of the shift, and no other part of the shift can be one.

\begin{proposition}[Commuting and anticommuting parts]
\label{bf:lp}
$\Ad_G:b\mapsto GbG^{-1}$ is an involution of $\so{3,3}$, so $\so{3,3}=\lG\oplus\mathfrak p_G$, where $\lG=\{b:[b,G]=0\}$ has dimension $6$ and
$\mathfrak p_G=\{b:bG=-Gb\}$ has dimension $9$. In the variable $z$,
\[
\lG=\{z\mapsto(P+iQ)z\}\cong\so{3,\C},\qquad \mathfrak p_G=\{z\mapsto(C+i\Sigma)\bar z\},
\]
with $P,Q,C$ real antisymmetric and $\Sigma$ real symmetric. The Killing form $4\Tr$ has signature $(3,3)$ on $\lG$. On $\mathfrak p_G$ it is
negative on the three directions $\Sigma=0$, compact counter-rotations of the two triples such as $R_{12}+R_{45}$, and positive on the six boosts
$C=0$. Both $\cK$, which acts as $z\mapsto-i\bar z$, and the pairing boosts $\sigma_j$ lie in $\mathfrak p_G$.
\end{proposition}

\begin{proof}
$G^{-1}=-G$ and $G^\dagger=G$ give $(GbG^{-1})^\dagger=-GbG^{-1}$, and $\Ad_G^2=\id$. In $(\bB,\bE)$ blocks, commuting with $G$ means the form
$\bigl(\begin{smallmatrix}P&-Q\\Q&P\end{smallmatrix}\bigr)$ and anticommuting means $\bigl(\begin{smallmatrix}C&\Sigma\\\Sigma&-C\end{smallmatrix}\bigr)$;
$b^\dagger=-b$ gives the symmetry types. The Killing form of $\so{3,3}$ is $4\Tr$, and on the two block forms $\Tr b^2=2\Tr(P^2-Q^2)$, respectively $2\Tr(C^2+\Sigma^2)$, which gives the signatures.
\end{proof}

This is fixed; the description of $\lG$ restates \cref{thm:GLorentz}. The pairing boosts that tear the binding at a threshold in \cref{prop:tear}
are anticommuting, and so is the grading root.

\paragraph{Gauge admissibility.} By \cref{thm:sig} and \cref{cor:noboost}, a Yang--Mills potential with a positive kinetic term takes values in
the compact part $\so3\oplus\so3$ of $\so{3,3}$, spanned by the $R_{ab}$ with $a,b$ in the same triple. On the two triples,
\begin{equation}
\label{bf:eq:compact}
b=\operatorname{diag}(A+D,\,A-D),\qquad A=[\mathbf a]_\times,\quad D=[\mathbf d]_\times,
\end{equation}
so the temporal triple rotates about $\mathbf a+\mathbf d$ and the spatial triple about $\mathbf a-\mathbf d$. The diagonal part $(A,A)$ lies in
$\lG$, with $w=\mathbf a$ real in \eqref{bf:eq:w}; the anti-diagonal part $(D,-D)$ lies in $\mathfrak p_G$, with $\Sigma=0$ (it acts as
$z\mapsto\mathbf d\times\bar z$). The boost-valued connections used below ($K_i$, $\sigma_j$, $\cK$, the $\Sigma$-part of $\mathfrak p_G$) are metric
connections in the sense of \cref{prop:tear}, legitimate in \eqref{bf:eq:L}, but they are not gauge fields. What a compact $b$ does pointwise to
the undivided energy and to the rest energy is \cref{en:transport}; here it acts through the interior dynamics.

\subsection{The phase-space generator and a criterion for boundedness}
\label{sec:bf-criterion}

\begin{lemma}[The phase-space generator]
\label{bf:M}
On $(X,\pi)$ the flow of \eqref{bf:eq:L} is generated by
\[
M=\begin{pmatrix}-S&\id/m\\0&S^\dagger\end{pmatrix},\qquad S^\dagger=-b+\kappa G,
\]
since $\dot X=\pi/m-SX$ and $\dot\pi=S^\dagger\pi$. It is the Hamiltonian flow of the Noether energy $H(X,p)=\tfrac1{2m}p^{\!\top}\eta^{-1}p-p^{\!\top}SX$:
in the canonical variables $(X,p)$ the generator is $\mathbb J\mathcal H$, with $\mathcal H$ the Hessian of $H$ and
$\mathbb J=\bigl(\begin{smallmatrix}0&\id\\-\id&0\end{smallmatrix}\bigr)$, and it is similar to $M$. $H$ vanishes on the invariant Lagrangian subspace
$p=0$, and $\det(\lambda-M)=\chi_S(\lambda)\chi_S(-\lambda)$ with $\chi_S(\lambda)=\det(\lambda-S)$, because $S^\dagger$ is similar to $S^{\!\top}$.
\end{lemma}

This is Hamiltonian bookkeeping, inherited. The classification below rests on the following criterion.

\begin{theorem}[When the interior motion stays bounded]
\label{bf:criterion}
Let $S$ be any constant real $6\times6$ matrix.
\begin{enumerate}[label=(\roman*),leftmargin=2.2em]
\item $M$ is semisimple if and only if $S$ is semisimple and, for every $\mu\in\C$, $\mathcal E_\mu(S^\dagger)$ and $\mathcal E_{-\mu}(S^\dagger)$ are
orthogonal for the bilinear extension of $\eta$.
\item Every solution is bounded on $\R$ if and only if, in addition, $\spec S\subset i\R$. Then (i) says that every eigenspace of $S^\dagger$ is
totally isotropic for $h_\eta$. Equivalently, every momentum normal mode $\pi=\Re(y\,e^{i\omega\tau})$ with $\omega\neq0$, superpositions within a
degenerate frequency included, has period mean $\langle\eta(\pi,\pi)\rangle=\half h_\eta(y,y)=0$, and the static momenta ($\omega=0$) span a totally
$\eta$-null space.
\end{enumerate}
\end{theorem}

\begin{proof}
$M$ is semisimple if and only if $-S$ (its restriction to $\pi=0$) and $S^\dagger$ (the induced map on the quotient) are semisimple and $M$ is
similar to $\operatorname{diag}(-S,S^\dagger)$, which by Roth's theorem \cite{Roth52} means that $-S\Xi-\Xi S^\dagger=-\id/m$ is solvable. With
$-S=F_1\Lambda_1F_1^{-1}$ and $S^\dagger=F_2\Lambda_2F_2^{-1}$ diagonal, this holds if and only if $(F_1^{-1}F_2)_{ij}=0$ whenever
$(\Lambda_1)_i=(\Lambda_2)_j=:\mu$. A row $r^{\!\top}$ of $F_1^{-1}$ with eigenvalue $\mu$ has $S^{\!\top}r=-\mu r$, and $S^{\!\top}=\eta S^\dagger\eta^{-1}$
puts $\eta^{-1}r$ in $\mathcal E_{-\mu}(S^\dagger)$; the condition is $r^{\!\top}y=\eta(\eta^{-1}r,y)=0$ for $y\in\mathcal E_\mu(S^\dagger)$. (ii) A linear
flow is bounded on $\R$ exactly when its generator is semisimple with imaginary spectrum, and then $\mathcal E_{-\mu}(S^\dagger)=\overline{\mathcal E_\mu(S^\dagger)}$.
\end{proof}

The theorem is fixed, although it holds for any nondegenerate $\eta$ and any $S$. Its necessity half is the split virial identity
(\cref{thm:splitvirial}: a bound orbit has null momentum on average); what it adds is that the averaged virial condition, imposed mode by mode, is
also sufficient. For $S=\kappa G$ it is the isotropy of \cref{prop:howmuch}.

\subsection{\texorpdfstring{Connections commuting with $G$: a Zeeman pattern without a scale}{Connections commuting with G: a Zeeman pattern without a scale}}
\label{sec:bf-commuting}

\begin{proposition}[The commuting part is a moving frame]
\label{bf:larmor}
Let $b\in\lG$. Then $\bigl(\tfrac d{d\tau}+b+\kappa G\bigr)(e^{-b\tau}Y)=e^{-b\tau}(\dot Y+\kappa GY)$. The exponents are therefore those of the
binding, $\pm i\kappa$, shifted by the eigenvalues of $-b$, and with $p_X=e^{b^{\!\top}\tau}p_Y$ the energies in the two frames are related by
\[
H_{\rm lab}=H_{\rm co}-p_Y^{\!\top}bY.
\]
\end{proposition}

\begin{proof}
$e^{-b\tau}$ is an isometry that commutes with $b$ and $G$; moreover $e^{b\tau}\eta^{-1}e^{b^{\!\top}\tau}=\eta^{-1}$ and
$e^{b\tau}(b+\kappa G)e^{-b\tau}=b+\kappa G$.
\end{proof}

For compact $b$ this is Larmor's theorem \cite{Larmor1897}, inherited: a rotating frame trades the field for an angular velocity. It is also the
first statement of \cref{prop:tear}. For a boost in $\lG$, $e^{-b\tau}$ is a uniformly boosted frame instead. $Q_b=p_Y^{\!\top}bY$ is the Noether
charge of the $\lG$ generator, so $H_{\rm lab}$ carries a chemical potential for it.

\begin{theorem}[Connections commuting with $G$]
\label{bf:lG}
Let $b\in\lG$ act as in \eqref{bf:eq:w}.
\begin{enumerate}[label=(\roman*),leftmargin=2.2em]
\item $\Xi=-G/(2\kappa m)$ solves $-S\Xi-\Xi S^\dagger=-\id/m$, so $M$ is semisimple if and only if $S$ is. The argument uses only $G^2=-\id$,
$G^\dagger=G$ and $[b,G]=0$.
\item $\chi_S(\lambda)=(\lambda^2+\kappa^2)\bigl((\lambda-i\kappa)^2+w\cdot w\bigr)\bigl((\lambda+i\kappa)^2+\overline{w\cdot w}\bigr)$.
\item In the frame in which $b$ is constant, every solution is bounded if and only if $e^{-b\tau}$ is bounded, that is, if and only if $w=0$ or
$w\cdot w>0$. The frequencies are then $\kappa$ and $|\kappa\pm s|$ with $s=\sqrt{w\cdot w}$. For $w\cdot w=0\neq w$ the zero-momentum flow grows
like $\tau^2$; for $w\cdot w\notin[0,\infty)$ it grows exponentially. For a gauge-admissible $b$, $w=\mathbf a$ is real and every solution is bounded.
\item For the rotation $b=\beta(R_{12}-R_{45})=-\beta L_3$ the spectrum of $M$ is $\pm i\kappa$, $\pm i(\kappa+\beta)$, $\pm i(\kappa-\beta)$, each
double. For the boost $b=\beta K_3$ the exponents are $\pm\beta\pm i\kappa$: complex rather than real, so the orbits spiral outward at the binding
frequency. The growth is the boost factor $e^{-\beta K_3\tau}$ of \cref{bf:larmor}, and $K_3$ is not a gauge generator (\cref{cor:noboost}).
\end{enumerate}
\end{theorem}

\begin{proof}
(i) With $[b,G]=0$: $-(b+\kappa G)\Xi-\Xi(-b+\kappa G)=\frac1{2\kappa m}(bG-Gb+2\kappa G^2)=-\frac1m$. (ii) On $(V,G)\cong\C^3$, $S$ is
$z\mapsto i\kappa z+w\times z$, and $(w\times)^3=-(w\cdot w)(w\times)$ gives $w\times$ the eigenvalues $0,\pm is$; the real characteristic polynomial
of a complex-linear map is its complex one times the conjugate. (iii) By \cref{bf:larmor}, $X=e^{-b\tau}Y$ with $Y$ a solution at $b=0$, bounded by
(i) and (ii); conversely the zero-momentum solutions are $e^{-b\tau}e^{-\kappa G\tau}X_0$ with $e^{\kappa G\tau}$ bounded. And $e^{-b\tau}$ is bounded
exactly when $w\times$ is semisimple with imaginary eigenvalues: $s$ real, and not $w\cdot w=0\neq w$, where $N=w\times$ has
$N^2\bar w=|w|^2w\neq0=N^3$. For real $w$, $w\times$ is real antisymmetric. (iv) Here $s=\pm\beta$, respectively $\pm i\beta$; in the first case $S$ is
semisimple and (i) applies.
\end{proof}

The theorem is fixed, but its content deflates. Criterion~(iii) is a property of the frame in which $b$ is constant, not of the binding. In the
co-moving frame of \cref{bf:larmor} the binding survives every commuting connection, as the remark after \cref{prop:tear} anticipates (growth on
one worldline proves nothing), and for admissible $b$ the motion in the frame in which $b$ is constant is bounded as well. In that frame, not in the co-moving one, a rotation lifts the degeneracy of the single
frequency across the three pairs, noted at the end of \cref{sec:neutral}, into the triplet $\kappa-\beta$, $\kappa$, $\kappa+\beta$. Read as a
frame vector through the identification of the frame with $\lG$ (\cref{thm:brackets}, with $L_i=e_i$ and $K_i=e_{i+3}$), the connection
has the invariant $q=w\cdot w$; on this reading the lab-frame motion stays bounded exactly for a connection of magnetic type, $\bE_b\cdot\bB_b=0$ and
$|\bB_b|>|\bE_b|$, while the binding itself survives every commuting connection, and an admissible one has $\bE_b=0$. The Zeeman pattern labels the
levels of the binding frequency by their $\lG$ charge. On one worldline it has no gauge-invariant content, and it is observable only against a
reference that $b$ does not rotate, as the Zeeman effect is observed against the laboratory; that is a reading, but a sharp one. A commuting $b$
supplies a spectrum and no invariant scale, since $\beta/\kappa$ is free.

\subsection{\texorpdfstring{Connections anticommuting with $G$, and screening}{Connections anticommuting with G, and screening}}
\label{sec:bf-anticommuting}

\begin{theorem}[Connections anticommuting with $G$]
\label{bf:pG}
Let $b\in\mathfrak p_G$ with $b\neq0$.
\begin{enumerate}[label=(\roman*),leftmargin=2.2em]
\item $(\pm b+\kappa G)^2=b^2-\kappa^2$, and $M^2=\operatorname{diag}(b^2-\kappa^2,\,b^2-\kappa^2)+N$ with
$N=\bigl(\begin{smallmatrix}0&-2b/m\\0&0\end{smallmatrix}\bigr)\neq0$. Hence $M$ is never semisimple, and some solutions grow at least linearly.
\item For the admissible counter-rotation $b=\beta(R_{12}+R_{45})$, which is $(D,-D)$ with $\mathbf d=-\beta\mathbf e_3$,
$\det(\lambda-M)=(\lambda^2+\kappa^2)^2(\lambda^2+\kappa^2+\beta^2)^4$. The exponents $\pm i\sqrt{\kappa^2+\beta^2}$ have algebraic multiplicity four and
geometric multiplicity two, and $\pm i\kappa$ are semisimple. The bounded solutions are exactly those with $\pi_1=\pi_2=\pi_4=\pi_5=0$, that is, with
no momentum on the support of $b$. The zero-momentum orbits are bounded and turn at the rate $\sqrt{\kappa^2+\beta^2}$ on the support.
\item The grading root as a connection satisfies $(\nu\cK+\kappa G)^2=(\nu^2-\kappa^2)\id$ on all of $V$.
\end{enumerate}
\end{theorem}

\begin{proof}
(i) $\{b,G\}=0$ gives both squares, and $M^2$ has off-diagonal block $(S^\dagger-S)/m=-2b/m$. With $b^2-\kappa^2=Q_s+Q_n$ (Jordan--Chevalley),
$M^2=\operatorname{diag}(Q_s,Q_s)+\bigl(\operatorname{diag}(Q_n,Q_n)+N\bigr)$ with commuting summands ($Q_n$ is a polynomial in $b$), so by uniqueness of the
decomposition $M^2$, and hence $M$, is semisimple only if $\operatorname{diag}(Q_n,Q_n)+N=0$, which is impossible for an off-diagonal $N\neq0$. A
non-semisimple generator of a Hamiltonian flow has a Jordan block at some exponent with non-negative real part. (ii) On the plane
$\langle e_3,e_6\rangle$, $b$ vanishes and $S=\kappa G$, semisimple with exponents $\pm i\kappa$, each double, by \cref{bf:lG}(i). On the support
$\langle e_1,e_2,e_4,e_5\rangle$, $b^2=-\beta^2$, so by (i) $S^2=(S^\dagger)^2=-\omega^2$ with $\omega^2=\kappa^2+\beta^2$ and, on the corresponding
eight-dimensional block of phase space, $M^2+\omega^2=N$ with $N^2=0$ and $\rank N=\rank b=4$. Hence the exponents there are $\pm i\omega$, each of
algebraic multiplicity four, $\ker(M^2+\omega^2)=\ker N$ has dimension four and splits into $\ker(M\mp i\omega)$ of dimension two each, and the bounded
data on the block are $\ker N=\{b\pi=0\}$. The zero-momentum flow is $e^{-S\tau}$ with $S^2=-\omega^2$ on the support. (iii) $\cK^2=\id$ and
$\{\cK,G\}=0$.
\end{proof}

Part~(i) is an obstruction: no purely anticommuting connection, admissible or not, is compatible with bounded normal modes. In particular the
Jordan blocks found after \cref{prop:tear} for a pairing-plane boost $\nu\sigma_j$ below the threshold occur for every $\nu\neq0$, below, at and above
it. Parts~(ii) and~(iii) are fixed. The rate changes to $\sqrt{\kappa^2+|\mathbf d|^2}$ for an admissible element and to $\sqrt{\kappa^2-\nu^2}$ for a
pairing boost (a metric connection, not a gauge field), and the momentum sector behaves as under the positive form of \cref{thm:virial}:
exponents carried by Jordan blocks, imaginary for a compact element and for a boost below the threshold.

\begin{proposition}[What a gauge transformation removes]
\label{bf:remove}
Let $b=b_l+b_p\in\lG\oplus\mathfrak p_G$ be constant. Some $g(\tau)\in O(3,3)$ carries $(b,G)$ to $(0,G')$ with $G'$ constant if and only if $[b,G]=0$.
For $b\in\mathfrak p_G$, $e^{b\tau}Ge^{-b\tau}=e^{2b\tau}G$. The gauge-covariant remainder is $\nabla_{\!\tau}G=[b,G]=2b_pG$: along a worldline $b_l$ is
pure gauge and $b_p$ is not.
\end{proposition}

\begin{proof}
$b^g=0$ forces $g=e^{-b\tau}g_0$, and then $G'=g_0^{-1}e^{b\tau}Ge^{-b\tau}g_0$ has derivative $g_0^{-1}e^{b\tau}[b,G]e^{-b\tau}g_0$. If $b$ anticommutes
with $G$, then $Gf(b)=f(-b)G$ for every power series $f$.
\end{proof}

This is fixed. A connection on a line has no curvature, since the non-abelian term $[b_\mu,b_\nu]$ needs two directions (inherited); once the
binding root is present, the invariant is $\nabla_{\!\tau}G$, which is exactly the anticommuting part. One might expect boundedness to force
$b\in\lG$. It does not, and the admissible case shows it.

\begin{theorem}[Compact screening]
\label{bf:screen}
Let $b$ be admissible, as in \eqref{bf:eq:compact}, and $\kappa>0$.
\begin{enumerate}[label=(\roman*),leftmargin=2.2em]
\item If every solution is bounded, then $\mathbf a\cdot\mathbf d=0$: the two triples rotate at equal rates, $|\mathbf a+\mathbf d|=|\mathbf a-\mathbf d|$.
\item Let $\mathbf a\cdot\mathbf d=0$ and $\mathbf a\neq0$. Then $\spec S=\{\pm i\omega_1,\pm i\omega_2,\pm i\omega_3\}$, where the $\omega_j$ are the
roots of
\[
f(\omega)=\omega^3+\kappa\omega^2-\bigl(\kappa^2+|\mathbf a|^2+|\mathbf d|^2\bigr)\,\omega-\kappa\bigl(\kappa^2+|\mathbf d|^2-|\mathbf a|^2\bigr),
\]
and if $|\mathbf a|^2\neq\kappa^2+|\mathbf d|^2$ every solution is bounded, although $[b,G]=2b_pG\neq0$ when $\mathbf d\neq0$.
\item If $\mathbf a\cdot\mathbf d=0$, $\mathbf d\neq0$ and $|\mathbf a|^2=\kappa^2+|\mathbf d|^2$, the vector
$X_0=\bigl(\mathbf n,(\mathbf a+\mathbf d)\times\mathbf n/\kappa\bigr)$ with $\mathbf n=\mathbf a\times\mathbf d$ is static, $SX_0=0$, with
$\eta(X_0,X_0)=-2|\mathbf d|^2|\mathbf n|^2/\kappa^2\neq0$, and some solutions drift linearly.
\end{enumerate}
For $\mathbf a=0$ and $\mathbf d\neq0$, \cref{bf:pG} applies.
\end{theorem}

\begin{proof}
$S=b+\kappa G$ is real antisymmetric, since $b$ and $G$ are; so $S$ and $S^\dagger=-\eta S\eta=-b+\kappa G$ are normal with imaginary spectrum, and
by \cref{bf:criterion} the motion is bounded if and only if $P_\mu\eta P_\mu=0$ for the orthogonal spectral projectors $P_\mu$ of $S^\dagger$.
(i) Then $\Tr\bigl(\eta(S^\dagger)^2\bigr)=\sum_\mu\mu^2\Tr(P_\mu\eta P_\mu)=0$. But $\{b,G\}$ is block-off-diagonal, $\Tr\eta=0$ and
$\Tr[\mathbf v]_\times^2=-2|\mathbf v|^2$, so $\Tr\bigl(\eta(S^\dagger)^2\bigr)=\Tr(A+D)^2-\Tr(A-D)^2=-8\,\mathbf a\cdot\mathbf d$.
(ii) Let $\rho$ be the half-turn about $\mathbf a$ and $\hat\rho=\operatorname{diag}(\rho,\rho)$. Since $\rho\mathbf a=\mathbf a$ and $\rho\mathbf d=-\mathbf d$,
$\hat\rho b\hat\rho=\operatorname{diag}(A-D,A+D)=G^{-1}bG$, so $G':=G\hat\rho$ commutes with $b$ and $G$, hence with $S$ and $S^\dagger$. Moreover
$G'^2=-\id$ and $G'^\dagger=G'$, so, as for $G$ in \cref{prop:howmuch}, each eigenspace $\mathcal V_\pm$ of $G'$ (eigenvalue $\pm i$) is totally
$h_\eta$-isotropic: $-i\,h_\eta(y,y')=h_\eta(G'y,y')=h_\eta(y,G'y')=i\,h_\eta(y,y')$ on $\mathcal V_+$. In the $(\bB,\bE)$ blocks
$\mathcal V_+=\{(\mathbf x,-i\rho\mathbf x)\}$, and $S$ acts there as $\mathbf x\mapsto(\mathbf a+\mathbf d)\times\mathbf x+i\kappa\rho\mathbf x$. After a
common rotation of both triples, which commutes with $G$ and $\eta$, to $\mathbf a=|\mathbf a|\mathbf e_3$ and $\mathbf d=|\mathbf d|\mathbf e_1$, one finds
$\det\bigl(i\omega-[\mathbf a+\mathbf d]_\times-i\kappa\rho\bigr)=-i f(\omega)$. So $S$ has the eigenvalues $i\omega_j$ on $\mathcal V_+$ and $-i\omega_j$ on
$\mathcal V_-=\overline{\mathcal V_+}$; since $\eta$ exchanges $\mathcal V_\pm$, $S^\dagger$ has the same eigenvalues there. Every eigenspace of $S^\dagger$
splits along $\mathcal V_\pm$, and lies in one of them unless $f(\omega)=f(-\omega)=0$ for some $\omega$. But
$\half[f(\omega)+f(-\omega)]=\kappa\bigl(\omega^2-\kappa^2-|\mathbf d|^2+|\mathbf a|^2\bigr)$ and
$\half[f(\omega)-f(-\omega)]=\omega\bigl(\omega^2-\kappa^2-|\mathbf a|^2-|\mathbf d|^2\bigr)$, so a common root forces $\mathbf a=0$, or $\omega=0$ and
$|\mathbf a|^2=\kappa^2+|\mathbf d|^2$. (iii) With $\mathbf u=\mathbf a+\mathbf d$ and $\mathbf v=\mathbf a-\mathbf d$,
$SX_0=\bigl(\mathbf u\times\mathbf n-\mathbf u\times\mathbf n,\ \mathbf v\times(\mathbf u\times\mathbf n)/\kappa+\kappa\mathbf n\bigr)=0$, because
$\mathbf v\cdot\mathbf n=0$ and $\mathbf u\cdot\mathbf v=|\mathbf a|^2-|\mathbf d|^2=\kappa^2$; and $\eta(X_0,X_0)=|\mathbf n|^2-|\mathbf u|^2|\mathbf n|^2/\kappa^2$.
So $\mathcal E_0(S^\dagger)=\eta\,\mathcal E_0(S)$ is not $\eta$-null, and \cref{bf:criterion} fails at $\mu=0$.
\end{proof}

The theorem is fixed, and it is the case that matters for a gauge field. Its mechanism is a second binding root $G'=G\hat\rho$ commuting with $S$. The exceptional ratio $|\mathbf a|^2=\kappa^2+|\mathbf d|^2$ has
measure zero. At $\mathbf d=0$ the theorem gives the Zeeman lines of \cref{bf:lG}. For $\mathbf d\neq0$ the screened frequencies cannot be removed
along a worldline (\cref{bf:remove}), but they are $\kappa$ times functions of $|\mathbf a|/\kappa$ and $|\mathbf d|/\kappa$: screening supplies no
scale.

\begin{proposition}[Gyroscopic screening by a boost]
\label{bf:gyro}
Let $b=\alpha L_3+b_\Sigma$, where $b_\Sigma\in\mathfrak p_G$ is the boost $z\mapsto i\Sigma\bar z$ with
$\Sigma=\bigl(\begin{smallmatrix}y_1&y_2&0\\y_2&-y_1&0\\0&0&0\end{smallmatrix}\bigr)$, and let $r^2=y_1^2+y_2^2>0$. Then $[b,G]\neq0$ and
\[
\chi_S(\lambda)=(\lambda^2+\kappa^2)\bigl(\lambda^2+(\kappa-\alpha)^2-r^2\bigr)\bigl(\lambda^2+(\kappa+\alpha)^2-r^2\bigr).
\]
If $\alpha\neq0$ and $r<\min(|\kappa-\alpha|,|\kappa+\alpha|)$, every solution is bounded; for $\alpha=0$ and for $|\kappa-\alpha|<r<|\kappa+\alpha|$ some
are not. Numerically, near bounded points at which the five conditions of the proof are independent, the set of $b\in\so{3,3}$ with all motion
bounded is a $10$-dimensional submanifold; at the tested points of $\lG$ and of this family the conditions are degenerate, so these points lie on
its singular locus.
\end{proposition}

\begin{proof}
On $\langle e_1,e_2,e_4,e_5\rangle$, $K_3=GL_3$ is an $\eta$-antisymmetric involution commuting with $G$, $L_3$ and $b_\Sigma$, so its eigenplanes
$\mathcal U_\pm$ are totally $\eta$-null and invariant. There $L_3=\mp G$, so $S=(\kappa\mp\alpha)G+b_\Sigma$ and $S^\dagger=(\kappa\pm\alpha)G-b_\Sigma$, each
squaring to a multiple of $\id$ ($b_\Sigma^2=r^2$, $\{b_\Sigma,G\}=0$), negative under the stated bound on $r$. For $\alpha\neq0$ the restrictions of
$S^\dagger$ to $\mathcal U_+$ and to $\mathcal U_-$ have no common eigenvalue, so every eigenspace of $S^\dagger$ lies in $\mathcal U_+\oplus\langle e_3,e_6\rangle$ or in
$\mathcal U_-\oplus\langle e_3,e_6\rangle$ (complexified). Since $\mathcal U_\pm$ are totally $\eta$-null and $\eta$-orthogonal to the $b$-free plane
$\langle e_3,e_6\rangle$, on which the eigenspaces of $\kappa G$ are isotropic lines, $h_\eta$ vanishes on every eigenspace, and \cref{bf:criterion}
applies. For $\alpha=0$, $b\in\mathfrak p_G$ and
\cref{bf:pG} applies; for $|\kappa-\alpha|<r<|\kappa+\alpha|$, $\chi_S$ has a real root. For the dimension statement: at simple frequencies
\cref{bf:criterion}(ii) imposes five real conditions on $b$ (two for imaginary spectrum after the trace relation, three for isotropy). Their
finite-difference Jacobian has rank five at four least-squares solutions and at an exact compact screened point, and rank four, with residuals of
second order along the kernel directions, at three points of this family and at two elliptic points of $\lG$.
\end{proof}

The family is fixed; the dimension statements are numerical. This is the gyroscopic stabilisation of a saddle
\cite{ThomsonTait1879,Kirillov2013}: the rotation detunes the two light-like planes of $K_3$ in which the boost acts, so that
$\eta(\pi,\pi)=2\eta(\pi_+,\pi_-)$, with $\pi_\pm$ the components of $\pi$ in the null planes $\mathcal U_\pm$ of the proof, averages to zero. Since $b_\Sigma$ is a boost, the family is not a gauge field. Together with \cref{bf:screen} it
shows that an anticommuting part need not cause secular growth, and that boundedness does not single out $\lG$; gauge removability does
(\cref{bf:remove}).

\begin{proposition}[A connection alone does not bind]
\label{bf:alone}
\begin{enumerate}[label=(\roman*),leftmargin=2.2em]
\item Let $\kappa=0$ and let $b\in\so{3,3}$ be any constant connection, $b=\nu\cK$ included. Then $X=e^{-b\tau}Y$ gives
$\mathcal L=\tfrac m2\eta(\dot Y,\dot Y)$ and $\ddot Y=0$. The generator $M=\bigl(\begin{smallmatrix}-b&\id/m\\0&-b\end{smallmatrix}\bigr)$ has
$\det(\lambda-M)=\chi_{-b}(\lambda)^2$ and is never semisimple: the motion is free drift, and there is no binding frequency.
\item The rescaling $(\tau,\kappa,b,m)\mapsto(\vartheta\tau,\kappa/\vartheta,b/\vartheta,\vartheta m)$ maps $M$ to $M/\vartheta$, solutions to solutions and
the action $\int\mathcal L\,d\tau$ to itself, and leaves $\eta$, $G$, $\lG$, $\mathfrak p_G$ and $\so4_\Delta$ of \cref{thm:PSbinding} unchanged. No value of $\kappa$, hence of
$\ell=c/\kappa$, is selected, with or without a constant $b$.
\item $M(m)=TM(1)T^{-1}$ with $T=\operatorname{diag}(\id,m\id)$: the spectrum of $M$ does not involve $m$, and the rescaling of (ii) keeps $m\kappa$,
not $m/\kappa$, fixed.
\end{enumerate}
\end{proposition}

\begin{proof}
(i) $\dot X+bX=e^{-b\tau}\dot Y$. Write $-\operatorname{diag}(b,b)=S'+N'$ (Jordan--Chevalley). Then $M=S'+(N'+N)$ with
$N=\bigl(\begin{smallmatrix}0&\id/m\\0&0\end{smallmatrix}\bigr)$, and $N'+N\neq0$ because $N$ is off-diagonal and $N'$ is block-diagonal. (ii) and (iii) are
direct.
\end{proof}

Part~(i) is inherited, since every connection on a line is pure gauge, and its consequence is an obstruction: a Yang--Mills potential on a
worldline, with no self-adjoint shift, produces free motion. Part~(ii) is dimensional analysis, inherited; it restates that $\kappa$ is not
fixed (end of \cref{sec:neutral}) and that no ratio among the inertia constants and $m$ is fixed by the Clifford structure (after
\cref{prop:several}). Part~(iii) is fixed and elementary: nothing in the interior dynamics ties $m$ to $\kappa$, and $mc^2=\hbar\kappa$, that is
$\ell=\hbar/(mc)$, is a matching added from outside, never a result.

\subsection{The split virial identity and the vacuum}
\label{sec:bf-vacuum}

The split virial identity of \cref{thm:splitvirial} holds unchanged with the connection: for every constant $S$, and so for every
$b\in\so{3,3}$, $\tfrac d{d\tau}\eta(X,\pi)=\eta(\pi/m-SX,\pi)+\eta(X,S^\dagger\pi)=\eta(\pi,\pi)/m$ along every solution. The quantum statement
needs the signature of the energy on each frequency.

\begin{theorem}[Krein signatures and the zero-point energy]
\label{bf:krein}
Suppose every solution is bounded, and let $\omega>0$ have multiplicity $k_\omega$ in $\spec(-iS)$. On the eigenspace $\mathcal E_{i\omega}$ of the
generator, the energy form $v\mapsto\bar v^{\!\top}\mathcal Hv$ has signature $(k_\omega,k_\omega)$. Hence, with positive-norm quantisation of the normal
modes, the zero-point energy $\sum_j(\pm)\half\hbar\omega_j$ cancels frequency by frequency, and every level has a lower one.
\end{theorem}

\begin{proof}
Work in the canonical variables of \cref{bf:M}. On $\mathcal E_{i\omega}$, $\mathcal Hv=-i\omega\mathbb Jv$, so the energy form equals $\omega$ times the
Krein form $-i\bar v^{\!\top}\mathbb Jv$, which is nondegenerate there because eigenspaces of distinct imaginary eigenvalues are orthogonal for it
\cite{YakubovichStarzhinskii75}. The zero-momentum part of $\mathcal E_{i\omega}$ is $\mathcal E_{i\omega}(-S)$, of dimension $k_\omega$, and $H$ vanishes on
it (\cref{bf:M}). Since the generator is semisimple, $\dim\mathcal E_{i\omega}=2k_\omega$ by \cref{bf:M}. A nondegenerate Hermitian form with a
half-dimensional isotropic subspace has signature $(k,k)$, and a mode of energy sign $\pm$ quantises to $\pm\hbar\omega(n+\half)$.
\end{proof}

The theorem is fixed, although structural: it holds for any nondegenerate $\eta$. It extends the vanishing of the zero-point energy on
$\R^{3,3}$ at $b=0$ (\cref{thm:subfs}) and the absence of a ground state (\cref{mb:noground}) to every connection that keeps the motion bounded: the
cancellation, forced by the perfect-square form \eqref{bf:eq:L}, survives the splitting line by line, and so does the absence of a ground state.
On the coordinate subfields of \cref{thm:subfs}, where the zero-point energy is $\pm\tfrac32\hbar\kappa$ or $\pm\half\hbar\kappa$, the theorem does not
apply, since $G$ does not preserve the subfield and the zero-momentum flow leaves it (\cref{thm:subfs}).

\begin{corollary}[No Fock vacuum with Jordan blocks]
\label{bf:nofock}
Suppose $M$ is not semisimple at a frequency $\omega\neq0$, as for the counter-rotation of \cref{bf:pG}(ii) and for a pairing boost below the threshold. Then no linear symplectic change of
variables writes the Hamiltonian on that block as a sum of signed oscillators, and there is no Fock vacuum annihilated by normal-mode lowering
operators.
\end{corollary}

\begin{proof}
A sum of signed oscillators generates a semisimple flow, and symplectic changes of variables act on the generator by similarity.
\end{proof}

This is an obstruction. A model block with the same Jordan structure is $H=\omega(q_1p_2-q_2p_1)+\half(q_1^2+q_2^2)$, whose two terms
Poisson-commute \cite{Williamson36}; that its quantisation has continuous spectrum is inherited and is not re-derived here.

\subsection{\texorpdfstring{Three constituents: the binding-compatible algebra $\so4_\Delta$}{Three constituents: the binding-compatible algebra so(4)Delta}}
\label{sec:bf-three}

At three constituents we use the generators $E_A$ of \eqref{eq:tower}, the numbering of \cref{sec:closure} (the temporal directions, their
partners $i'=i+3$, the weak directions $W=\{7,13,14,15\}$, the colour directions $\{10,11,12,16,17,18\}$ and the span $E$), $G_3=G\oplus G\oplus G$,
and $\mathrm{rot}(i,j)$ acting by $\tfrac12\ad$; $R_{AB}$ is defined as at one constituent. By \cref{thm:PSbinding}, the part of the Pati--Salam
algebra of $P_*$ that commutes with $G_3$ is $\so4_\Delta=\operatorname{span}\{\mathrm{rot}(i,j):i<j\in W\}\cong\su2\oplus\su2$, compact, and neither
Pati--Salam factor meets $\Cent(G_3)$. The spectra below use an explicit splitting.

\begin{lemma}[The two factors of $\so4_\Delta$]
\label{bf:su2}
The elements $T^\pm_1=\half\bigl(\mathrm{rot}(13,14)\pm\mathrm{rot}(7,15)\bigr)$, $T^\pm_2=\half\bigl(\mathrm{rot}(7,14)\mp\mathrm{rot}(13,15)\bigr)$ and
$T^\pm_3=\half\bigl(\mathrm{rot}(7,13)\pm\mathrm{rot}(14,15)\bigr)$ satisfy $[T^+_i,T^-_j]=0$ and $[T^\pm_1,T^\pm_2]=-T^\pm_3$ cyclically, and the trace
form is negative definite on their span, so $\so4_\Delta=\su2_+\oplus\su2_-$. On the eighteen vectors each $T^\pm_3$ has the eigenvalues $\pm\tfrac i2$,
four times each, and $0$, ten times.
\end{lemma}

\begin{proof}
Each $\mathrm{rot}(i,j)$ with $i,j\in W$ rotates the plane $(e_i,e_j)$ and its partner plane $(e_{i'},e_{j'})$ alike, so on $\operatorname{span}W$, and
identically on the partner span $\{10,16,17,18\}$, the six elements act as the standard generators of $\so4$, and the $T^\pm_i$ are their
self-dual and anti-self-dual combinations; the brackets follow by direct computation. Each $T^\pm_3$ is half a sum of two commuting plane rotations
on each of the two four-dimensional spans and vanishes on the remaining ten directions.
\end{proof}

The lemma is fixed; it makes the isomorphism $\so4_\Delta\cong\su2\oplus\su2$ of \cref{thm:PSbinding} explicit. Each $\su2$ factor locks a weak rotation of $W$ to the colour rotation of the partners $\{10,16,17,18\}$. Reading $\so4_\Delta$ as Yang and Mills'
isospin algebra, doubled, is a reading, as is the description after \cref{thm:survive}.

\begin{proposition}[Spectra on the eighteen-dimensional frame]
\label{bf:so4spec}
Let $\kappa=5$ and $m=1$.
\begin{enumerate}[label=(\roman*),leftmargin=2.2em]
\item For $b=\beta T^+_3$ with $\beta=2$, $\chi_S=(\lambda^2+25)^5(\lambda^2+16)^2(\lambda^2+36)^2$: the frequency $\kappa$ on the five complex singlet
directions, and $\kappa\mp\beta/2$ on the two $\su2_+$ doublets, of weights $\pm\half$.
\item For $b=\beta_+T^+_3+\beta_-T^-_3$ with $(\beta_+,\beta_-)=(2,4)$, the lines are $\kappa\pm\half(\beta_+\pm\beta_-)=8,6,4,2$, together with $\kappa$.
\item In both cases $\Xi=-G_3/(2\kappa m)$ solves the Sylvester equation, $M$ is semisimple, the Krein signatures are $(k,k)$, and the zero-point
energy vanishes.
\item The weak-only rotation $R_{7,13}$ and the colour-only rotation $R_{10,16}$ give a non-semisimple $M$, hence secular growth.
\item $b=R_{7,13}\pm R_{10,11}\in\mathfrak{ps}_W$ has $[b,G_3]\neq0$, and every solution is bounded, for every $\kappa>0$. At $\kappa=5$ the frequencies
are $\kappa$ (six pairs) and $|\omega_j|\approx4.319,\,5.050,\,5.731$, with $\omega_j$ the roots of $\omega^3+5\omega^2-26\omega-125$.
\end{enumerate}
\end{proposition}

\begin{proof}
(i)--(iii) The connection lies in $\so4_\Delta\subset\Cent(G_3)$ and is real antisymmetric, so \cref{bf:lG}(i), whose argument uses only
$G_3^2=-\id$, $G_3^\dagger=G_3$ and $[b,G_3]=0$, gives semisimplicity and boundedness, and \cref{bf:krein} gives the signatures and the zero-point
energy. Since $b$ commutes with $G_3$, $S=b+\kappa G_3$ is complex linear on $(\R^{18},G_3)\cong\C^9$ and acts as $i\kappa+b$. By the proof of \cref{bf:su2}, on the four
complex directions of $W$ the commuting $T^+_3,T^-_3$ have the joint eigenvalues $\pm\tfrac i2(1,1)$ and $\pm\tfrac i2(1,-1)$, and both vanish on the other five; the
characteristic polynomials follow as in the proof of \cref{bf:lG}(ii). (iv) and (v) $G_3$ maps the six-plane $\langle e_7,e_8,e_{13};e_{10},e_{11},e_{16}\rangle$ as
$G$ maps $\langle e_1,e_2,e_3;e_4,e_5,e_6\rangle$, and each $b$ named preserves that six-plane and vanishes on its $\eta$-orthogonal complement, where
the motion is the binding at $b=0$. On the six-plane $b$ is admissible. $R_{7,13}$ rotates the temporal triple alone and $R_{10,16}$ the spatial
triple alone, so $\mathbf a\cdot\mathbf d=\pm|\mathbf a|^2\neq0$, and \cref{bf:screen}(i) applies; since $S$ is real antisymmetric, the unbounded
solutions grow polynomially. For $R_{7,13}\pm R_{10,11}$ the rotation vectors of the two triples are the unit vectors $\mathbf e_2$ and $\pm\mathbf e_3$:
equal rates and perpendicular axes, so $\mathbf a=\half(\mathbf e_2\pm\mathbf e_3)\perp\mathbf d=\half(\mathbf e_2\mp\mathbf e_3)$ and
$|\mathbf a|^2=|\mathbf d|^2=\half\neq\kappa^2+\half$. \Cref{bf:screen}(ii) applies, with $f$ as stated at $\kappa=5$; the three frequencies are its
roots, evaluated numerically.
\end{proof}

This is fixed. Every constant $\so4_\Delta$-valued connection is compact and commutes with $G_3$, so it keeps the interior bounded; the splitting
follows the $\su2$ weights, $\beta/\kappa$ is free, and by \cref{bf:larmor} the pattern can be removed along one worldline. The single-factor
rotations tear the binding, but boundedness does not single out the locked diagonal: the connection of part~(v) is bounded with $[b,G_3]\neq0$,
while the anti-locked $R_{7,13}+R_{10,16}$, which has $\mathbf a=0$ on its six-plane, is not. What singles out $\so4_\Delta$ is commutation with
$G_3$ (\cref{thm:PSbinding}), that is, removability along a worldline (\cref{bf:remove}).

\subsection{Quantum Yang--Mills facts, and the channels of the vacuum}
\label{sec:bf-quantum}

Classically the connection splits, screens or tears the binding but supplies no scale. The only place left for it to generate one is the
quantum theory, and the following inherited facts are needed to say what that would involve.

\begin{proposition}[Inherited facts]
\label{bf:qym}
\begin{enumerate}[label=(\roman*),leftmargin=2.2em]
\item For $\su4$ with $\tr(T^AT^B)=\half\delta^{AB}$, $\sum_AT^AZT^A=\half\bigl(\tr(Z)\id-\tfrac14Z\bigr)$ for every $4\times4$ matrix $Z$: the factor is
$\tfrac{15}8=C_F$ on $\id$ and $-\tfrac18$ on traceless $Z$.
\item The channel factors $T_1\cdot T_2=\half[C_2(\text{pair})-C_2(R_1)-C_2(R_2)]$ are $+\tfrac38,-\tfrac58$ on $\mathbf{10},\mathbf6\subset\mathbf4\otimes\mathbf4$,
$-\tfrac{15}8,+\tfrac18$ on $\mathbf1,\mathbf{15}\subset\bar{\mathbf4}\otimes\mathbf4$, and $+\tfrac14,-\tfrac34$ on $\mathbf3,\mathbf1\subset\mathbf2\otimes\mathbf2$ of
$\su2$, with $C_2=\tfrac{15}8,\tfrac92,\tfrac52,4$ on $\mathbf4,\mathbf{10},\mathbf6,\mathbf{15}$. Hence the Dirac bilinear
$\overline{(\mathbf4,\mathbf2,\mathbf1)}\otimes(\mathbf4,\mathbf1,\mathbf2)=(\mathbf1,\mathbf2,\mathbf2)\oplus(\mathbf{15},\mathbf2,\mathbf2)$ has $\su4$ factors
$-\tfrac{15}8$ and $+\tfrac18$, and a Lorentz-scalar Majorana bilinear of $(\mathbf4,\mathbf1,\mathbf2)$, symmetric in its internal indices, lies in
$\mathrm{Sym}^2(\C^4\otimes\C^2)=(\mathbf{10},\mathbf1,\mathbf3)\oplus(\mathbf6,\mathbf1,\mathbf1)$, of dimensions $30+6$, with factors
$(+\tfrac38,+\tfrac14)$ and $(-\tfrac58,-\tfrac34)$.
\item In Euclidean momentum space
$\int_{k^2\le k_{\max}^2}\frac{d^4k}{(2\pi)^4}\frac1{k^2+M_{\rm dyn}^2}=\frac1{16\pi^2}\bigl[k_{\max}^2-M_{\rm dyn}^2\ln(1+k_{\max}^2/M_{\rm dyn}^2)\bigr]$. The function
$f_{\rm gap}(x)=1-x^2\ln(1+x^{-2})$ decreases strictly from $1$ to $0$, so the contact gap equation $1=\lambda_{\rm c}f_{\rm gap}(x)$ has a root
$x=M_{\rm dyn}/k_{\max}>0$, with $M_{\rm dyn}$ the dynamical mass, $k_{\max}$ the cutoff and $\lambda_{\rm c}$ the dimensionless contact coupling, if and only if $\lambda_{\rm c}>1$.
\item The ladder critical coupling $\pi/(3C_F)$ \cite{FukudaKugo76,Miransky93} is $8\pi/45\approx0.5585$ at $C_F=\tfrac{15}8$. For $SU(2)$ with $N_f$
Dirac doublets the one-loop coefficient is $\beta_0=(22-2N_f)/3$, equal to $6$ at $N_f=2$ \cite{GrossWilczek73,Politzer73}, and
$\Lambda=\mu\,e^{-8\pi^2/(\beta_0g^2(\mu))}$ is independent of $\mu$. The breaking $SU(4)\to Sp(4)$ of $N_f=2$ leaves $15-10=5$ Goldstone
directions \cite{Peskin80}.
\end{enumerate}
\end{proposition}

All of this is inherited \cite{NambuJonaLasinio61,Slansky}; the strict decrease in (iii) follows from $\ln(1+u)>u/(1+u)$. The use of (ii) is a reading.
Under the most-attractive-channel heuristic of one-boson exchange \cite{RabyDimopoulosSusskind80}, in which a negative factor means attraction, and
with a common coupling for $\su4$ and $\su2_R$, which this paper adopts only as an assumption (\cref{ass:field} and the remark after
\cref{prop:coupled3}), the exchange favours the Dirac singlet channel $(\mathbf1,\mathbf2,\mathbf2)$ and the colour-breaking $(\mathbf6,\mathbf1,\mathbf1)$.
It disfavours the $(\mathbf{15},\mathbf2,\mathbf2)$ that the Georgi--Jarlskog factor $-3$ needs (\cref{prop:EW}; \cite{GeorgiJarlskog79}), and the
$(\mathbf{10},\mathbf1,\mathbf3)$ that carries the Standard-Model vacuum of \cref{sec:selection}. That vacuum is a minimum of a scalar potential, not a
gauge-exchange condensate, so the heuristic addresses a different, dynamical scenario, in which a gauge-generated vacuum would point away from the
one selected there. In the selected vacuum the order is moreover the reverse of a mass generated by the gauge field: the Pati--Salam vector bosons
outside the Standard Model acquire their masses by eating nine Goldstone directions (\cref{thm:nomajoron}).

\subsection{What the connection part supplies}
\label{sec:bf-synthesis}

The results of this section combine as follows.
\begin{enumerate}[label=(\alph*),leftmargin=2.2em]
\item The binding frequency comes from the self-adjoint part of the shift, never from a connection. The rate $\kappa$, with
$(\kappa G)^2=-c^2/\ell^2$, sits in the tensor part (\cref{bf:covariant}, fixed), and a connection alone gives free drift (\cref{bf:alone}, an
obstruction).
\item A commuting connection splits the binding frequency without supplying a scale: the lines $\kappa$ and $\kappa\pm\sqrt{w\cdot w}$ are bounded for
admissible $b$ and removable along any worldline, with $\beta/\kappa$ free (\cref{bf:lG,bf:larmor}, fixed; the Zeeman description is a reading). At
three constituents $\so4_\Delta$ gives lines labelled by $\su2$ weights and is not the only bounded part of the Pati--Salam algebra
(\cref{bf:so4spec}), and the zero-point cancellation survives line by line (\cref{bf:krein}).
\item A purely anticommuting connection tears the binding (\cref{bf:pG,bf:nofock}, obstructions), with rates $\sqrt{\kappa^2+|\mathbf d|^2}$ for an
admissible element or $\sqrt{\kappa^2-\nu^2}$ for a pairing boost, carried by Jordan blocks. A mixed admissible connection keeps every orbit bounded
exactly at equal rates of the two triples about axes that are not opposite, outside the ratio $|\mathbf a|^2=\kappa^2+|\mathbf d|^2$
(\cref{bf:screen}); a boost can be screened gyroscopically (\cref{bf:gyro}); neither gives a scale.
\item The only role left for the connection in generating a scale is quantum, and it is open. The arena leaves $\ell$ free (\cref{bf:alone}, and
\cref{sec:restclosure} for a dynamical rate), and classical four-dimensional Yang--Mills theory has no dimensionful parameter; a scale arises only
by dimensional transmutation, $\Lambda_b=\mu\,e^{-8\pi^2/(\beta_0g^2(\mu))}$ \cite{ColemanWeinberg73,GrossWilczek73}, for a Yang--Mills theory whose
potential is $b$.
\end{enumerate}

\begin{conjecture}[A transmuted binding length]
\label{bf:transmute}
$\ell=\hbar c/\Lambda_b$ for the scale $\Lambda_b$ of a Yang--Mills theory whose potential is the connection part of the shift; that is,
$\hbar\kappa=\Lambda_b$.
\end{conjecture}

The conjecture is open: nothing in the arena derives it. Reading the rest energy $mc^2$ as $\hbar\kappa$ is a separate identification and not part of
the conjecture, since the spectrum of $M$ does not involve $m$ and the rescaling of \cref{bf:alone} keeps $m\kappa$, not $m/\kappa$, fixed. A
mechanism would need $\kappa$ to be dynamical, as the rotor rate of \cref{sec:closure} is, and coupled to a quantum gauge sector with a fixed
$\beta_0$ and a matching coupling; even then, by \cref{bf:qym}(ii), the most attractive channels are not those of the vacuum selected in
\cref{sec:selection}. A classification of the compact connections that keep the binding bounded at three constituents, beyond the commuting ones and the equal-rate
family embedded as in \cref{bf:so4spec}(v), is also open, as is a description of the bounded set in $\so{3,3}$ beyond the numerical statement of
\cref{bf:gyro}.

Read in the variables of $E=mc^2$, and only through the added identification $mc^2=\hbar\kappa$, the results of this section complete the reading
of \cref{sec:energy-readings}, in which the self-adjoint summand $\kappa G$ carries $c$ and alone binds; what follows is a reading of the fixed
results above. The Yang--Mills potential is the antisymmetric, transport summand of the same shift, and it is compact. Classically it generates neither $m$ nor $\kappa$: when it commutes with $G$ it splits the interior rate $\kappa$
into charge-labelled lines, a purely anticommuting potential tears the binding, and a mixed one may or may not (\cref{bf:screen}). Whether
its quantum scale is the length $\ell$ that the arena leaves free is open.
